\section{Purpose, Scope, and Project Evolution}

\subsection{What the project is trying to establish}

A passive limit order does not receive a random sample of future price moves.
It fills when opposing flow consumes the modeled quantity ahead of it, while
other orders may remain unfilled or be cancelled. A signal can therefore
predict the market's next mid-price move without improving the outcome of the
particular trades a passive maker receives. The project asks whether that
distinction is visible in a controlled BTCUSDT experiment.

The deliverable is an offline chain from recorded data to a documented research
decision: capture and integrity inventory; sequence-valid book reconstruction;
deterministic replay; explicit execution and inventory constraints; economic
and signal diagnostics; a predeclared development gate; and verified artifacts.
Native implementation follows a behavioral contract so that speed can be
assessed after correctness. A profitable strategy is an empirical possibility,
not a requirement imposed on the answer.

This is a small-order counterfactual simulator over aggregate L2. It is not a
market-by-order reconstruction, a calibrated exchange replica, or a live
trading service. The distinction matters because an exact implementation of a
queue assumption does not make that assumption observable in the underlying
data.

\subsection{The bar-backtesting precursor}

The initial framework used BTCUSDT hourly OHLCV data and a common interface for
momentum, mean reversion, moving-average crossover, and Bollinger-band
strategies. Data loading, transaction costs, portfolio accounting, metrics,
parameter sweeps, and walk-forward orchestration were separate components.
Useful lessons included recreating stateful strategies between runs and
checking whether fees erased a gross result.

The historical Phase 1 ratios are not current strategy evidence. The audit
found same-close signal observation and execution, hourly equity annualized as
365 periods per year, an ``out-of-sample Sharpe'' formed from monthly Sharpe
estimates, and an optimizer source overwritten during plotting. The metrics
helper now infers observation frequency, and the logged search grids have a
reconstructed runner; the old CSVs retain their historical identity. These
limitations motivated the move to observable book events rather than a more
elaborate interpretation of coarse bar returns.

\subsection{Why the project moved to L2}

Bars cannot identify post-only eligibility at arrival, displayed volume ahead,
partial fills, cancellation races, or whether a fill is followed by an adverse
move. L2 depth and trades expose enough information to model these mechanisms,
while still leaving true queue position latent. The engineering design makes
that missing information an explicit parameter, rather than silently treating
a price touch as a guaranteed fill.

Initial L2 exploration rejected attractive results contaminated by unintended
taker execution. A later profitable passive anchor did not generalize across
the expanded panel. The fixed microprice baseline and OFI premise test then
became a structured development study, with a reserved holdout and a stopping
rule when the premise failed.

\begin{table}[H]
\centering
\caption{The project record preserves distinct research and implementation stages.}
\begin{tabularx}{\textwidth}{@{}L{0.16\textwidth}L{0.26\textwidth}Y@{}}
\toprule
Stage & Deliverable & Evidential role \\
\midrule
Jan--Apr 2026 & Bar framework and walk-forward exercise & Educational precursor; historical performance limitations are retained. \\
Apr--May 2026 & L2 capture, replay, baseline exploration & Established passive execution mechanics and the development hypothesis. \\
Jun 2026 & Frozen V2 24-window study & Historical baseline, OFI gate, queue-credit and latency stress under \LegacyModel. \\
Jul--Aug 2026 & Execution and publication audits & Exact private scheduling, own FIFO, working exposure, snapshot timing, and explicit model boundaries. \\
Sep 2026 & Corrected V3 development rerun & Same development data and locked baseline; current-model economics and diagnostic conditional uncertainty. \\
Sep 2026 & Native implementation and report closure & Checked C++ book, Python parity, scoped measurement, and this complete project account. \\
\bottomrule
\end{tabularx}
\end{table}

\subsection{Model corrections and the meaning of completion}

The frozen V2 engine activated orders at the next book update after their
modeled arrival and made cancellations effective immediately. A later audit
also identified snapshots stamped before a blocking REST request completed.
These details could change fill eligibility and queue age. The project froze
the old evidence as \LegacyModel, implemented \CurrentModel, and reran the
development experiment in a distinct V3 namespace.

The V3 protocol was committed at \ProtocolCommit{} before its corrected
economics were inspected; \ResearchCommit{} identifies the research source.
Completion means executing that protocol, reporting its actual outcome,
preserving the holdout boundary, and validating the native work within its
stated arithmetic and integration contract. It does not require advancing a
candidate that failed the gate or relabeling the historical sensitivity grid
as current-model evidence.
